% GENERATED FILE. Edit canonical lessons, then run npm run generate:books.
% canonical-source-hash: fnv1a64:3fe8b7bf79dfdb7d
% canonical-lessons: AR-C304-zaman, AR-C304-mundhu-zaman-tawil, AR-C304-hiqba, AR-C304-anadhak, AR-C304-ahd

\chapter{Long ago}
\label{ch:ar-long-ago}
\hlchaptermodality{304}

\begin{chapteropening}
\textbf{By the end of this chapter:} \emph{I can say an age, long ago, an era, at that time and a reign.}

Complete the last lesson of chapter 304: I can say an age, long ago, an era, at that time and a reign.
\end{chapteropening}

\section[\texorpdfstring{\ar{زمن} (zaman)}{زمن (zaman)}]{\ar{زمن} (zaman) — an age, a span of time}

\label{lesson:AR-C304-zaman}



\begin{quote}
Before the new one: say the Arabic for a villain, then the Arabic for to imagine.
\end{quote}

\subsection*{You'll want to know: \ar{زمن}}
\textbf{\ar{زمن}} — \emph{zaman} — \textquotedblleft{}an age, a span of time\textquotedblright{}.

\begin{quote}
\ar{كان} \ar{ذلك} \ar{في} \ar{زمن} \ar{جدي}.
\end{quote}

\begin{quote}
\emph{kāna dhālika fī zaman jaddī.}
\end{quote}

\begin{quote}
That was in my grandfather’s time.
\end{quote}

Plural: \ar{أزمنة}, \emph{azmina}.

\ar{وقت} is a time for something; \ar{زمن} is time as a stretch, an age.

\begin{cousinweb}
\textbf{\ar{زمن}} (\emph{zaman}) is built on the native Arabic root \emph{z-m-n} (Lane), where \emph{zaman} and \emph{zamān} name time, a short span or a long one.
\end{cousinweb}

\begin{grammarlens}[title={What you've built}]
One more word for telling what happened.
\end{grammarlens}

\subsection*{Your turn}
\begin{itemize}
\raggedright
  \item \emph{Say it:} \emph{zaman}
  \item \emph{Say it:} \emph{zaman}, once more
  \item \emph{Say it:} \emph{kāna dhālika fī zaman jaddī.} — That was in my grandfather’s time.
  \item \emph{Recall:} say the Arabic for a villain, then the Arabic for to imagine, then say \emph{zaman} again
\end{itemize}

\begin{tcolorbox}[breakable,colback=teal!4,colframe=teal!35!black,title={Before you move on}]
What does \ar{زمن} mean? (\textquotedblleft{}An age, a span of time\textquotedblright{}.) Say it once more.
\end{tcolorbox}
\section[\texorpdfstring{\ar{منذ} \ar{زمن} \ar{طويل} (mundhu zaman ṭawīl)}{منذ زمن طويل (mundhu zaman ṭawīl)}]{\ar{منذ} \ar{زمن} \ar{طويل} (mundhu zaman ṭawīl) — long ago; for a long time}

\label{lesson:AR-C304-mundhu-zaman-tawil}



\begin{quote}
Before the new one: say the Arabic for to imagine, then the Arabic for an age.
\end{quote}

\subsection*{You'll want to know: \ar{منذ} \ar{زمن} \ar{طويل}}
\textbf{\ar{منذ} \ar{زمن} \ar{طويل}} — \emph{mundhu zaman ṭawīl} — \textquotedblleft{}long ago; for a long time\textquotedblright{}.

\begin{quote}
\ar{منذ} \ar{زمن} \ar{طويل} \ar{كانت} \ar{هنا} \ar{غابة} \ar{كبيرة}.
\end{quote}

\begin{quote}
\emph{mundhu zaman ṭawīl kānat hunā ghāba kabīra.}
\end{quote}

\begin{quote}
Long ago there was a big forest here.
\end{quote}

It opens a story set in the past, or says how long a thing has lasted.

\begin{cousinweb}
\textbf{\ar{منذ} \ar{زمن} \ar{طويل}} (\emph{mundhu zaman ṭawīl}) is \textquotedblleft{}since a long time\textquotedblright{}: the taught \emph{mundhu}, \textquotedblleft{}since\textquotedblright{}, with \emph{zaman}, built on the native Arabic root \emph{z-m-n} (Lane), and \emph{ṭawīl}, long.
\end{cousinweb}

\begin{grammarlens}[title={What you've built}]
One more word for telling what happened.
\end{grammarlens}

\subsection*{Your turn}
\begin{itemize}
\raggedright
  \item \emph{Say it:} \emph{mundhu zaman ṭawīl}
  \item \emph{Say it:} \emph{mundhu zaman ṭawīl}, once more
  \item \emph{Say it:} \emph{mundhu zaman ṭawīl kānat hunā ghāba kabīra.} — Long ago there was a big forest here.
  \item \emph{Recall:} say the Arabic for to imagine, then the Arabic for an age, then say \emph{mundhu zaman ṭawīl} again
\end{itemize}

\begin{tcolorbox}[breakable,colback=teal!4,colframe=teal!35!black,title={Before you move on}]
What does \ar{منذ} \ar{زمن} \ar{طويل} mean? (\textquotedblleft{}Long ago; for a long time\textquotedblright{}.) Say it once more.
\end{tcolorbox}
\section[\texorpdfstring{\ar{حقبة} (ḥiqba)}{حقبة (ḥiqba)}]{\ar{حقبة} (ḥiqba) — an era, a long period}

\label{lesson:AR-C304-hiqba}



\begin{quote}
Before the new one: say the Arabic for an age, then the Arabic for long ago.
\end{quote}

\subsection*{You'll want to know: \ar{حقبة}}
\textbf{\ar{حقبة}} — \emph{ḥiqba} — \textquotedblleft{}an era, a long period\textquotedblright{}.

\begin{quote}
\ar{كانت} \ar{تلك} \ar{حقبة} \ar{صعبة} \ar{في} \ar{تاريخ} \ar{البلد}.
\end{quote}

\begin{quote}
\emph{kānat tilka ḥiqba ṣaʿba fī tārīkh al-balad.}
\end{quote}

\begin{quote}
That was a hard era in the country’s history.
\end{quote}

Plural: \ar{حقب}, \emph{ḥiqab}.

\begin{cousinweb}
\textbf{\ar{حقبة}} (\emph{ḥiqba}) is built on the native Arabic root \emph{ḥ-q-b} (Lane), where \emph{ḥiqba} and \emph{ḥuqb} name a long span of years, and \emph{ḥaqība}, a bag tied behind the saddle, is the same root.
\end{cousinweb}

\begin{grammarlens}[title={What you've built}]
One more word for telling what happened.
\end{grammarlens}

\subsection*{Your turn}
\begin{itemize}
\raggedright
  \item \emph{Say it:} \emph{ḥiqba}
  \item \emph{Say it:} \emph{ḥiqba}, once more
  \item \emph{Say it:} \emph{kānat tilka ḥiqba ṣaʿba fī tārīkh al-balad.} — That was a hard era in the country’s history.
  \item \emph{Recall:} say the Arabic for an age, then the Arabic for long ago, then say \emph{ḥiqba} again
\end{itemize}

\begin{tcolorbox}[breakable,colback=teal!4,colframe=teal!35!black,title={Before you move on}]
What does \ar{حقبة} mean? (\textquotedblleft{}An era, a long period\textquotedblright{}.) Say it once more.
\end{tcolorbox}
\section[\texorpdfstring{\ar{آنذاك} (ānadhāk)}{آنذاك (ānadhāk)}]{\ar{آنذاك} (ānadhāk) — at that time, back then}

\label{lesson:AR-C304-anadhak}



\begin{quote}
Before the new one: say the Arabic for long ago, then the Arabic for an era.
\end{quote}

\subsection*{You'll want to know: \ar{آنذاك}}
\textbf{\ar{آنذاك}} — \emph{ānadhāk} — \textquotedblleft{}at that time, back then\textquotedblright{}.

\begin{quote}
\ar{كان} \ar{أبي} \ar{آنذاك} \ar{طالبا} \ar{في} \ar{الجامعة}.
\end{quote}

\begin{quote}
\emph{kāna abī ānadhāk ṭāliban fī al-jāmiʿa.}
\end{quote}

\begin{quote}
My father was a university student at that time.
\end{quote}

\begin{cousinweb}
\textbf{\ar{آنذاك}} (\emph{ānadhāk}) is usually traced to \emph{ān}, \textquotedblleft{}a time\textquotedblright{}, the same word as in the taught \ar{الآن}, now, joined to \emph{dhāka}, \textquotedblleft{}that\textquotedblright{}: \textquotedblleft{}at that time\textquotedblright{}.
\end{cousinweb}

\begin{grammarlens}[title={What you've built}]
One more word for telling what happened.
\end{grammarlens}

\subsection*{Your turn}
\begin{itemize}
\raggedright
  \item \emph{Say it:} \emph{ānadhāk}
  \item \emph{Say it:} \emph{ānadhāk}, once more
  \item \emph{Say it:} \emph{kāna abī ānadhāk ṭāliban fī al-jāmiʿa.} — My father was a university student at that time.
  \item \emph{Recall:} say the Arabic for long ago, then the Arabic for an era, then say \emph{ānadhāk} again
\end{itemize}

\begin{tcolorbox}[breakable,colback=teal!4,colframe=teal!35!black,title={Before you move on}]
What does \ar{آنذاك} mean? (\textquotedblleft{}At that time, back then\textquotedblright{}.) Say it once more.
\end{tcolorbox}
\section[\texorpdfstring{\ar{عهد} (ʿahd)}{عهد (ʿahd)}]{\ar{عهد} (ʿahd) — a reign, an era; a pledge}

\label{lesson:AR-C304-ahd}



\begin{quote}
Before the new one: say the Arabic for an era, then the Arabic for at that time.
\end{quote}

\subsection*{You'll want to know: \ar{عهد}}
\textbf{\ar{عهد}} — \emph{ʿahd} — \textquotedblleft{}a reign, an era; a pledge\textquotedblright{}.

\begin{quote}
\ar{بني} \ar{الجسر} \ar{في} \ar{عهد} \ar{الملك} \ar{الأول}.
\end{quote}

\begin{quote}
\emph{buniya al-jisr fī ʿahd al-malik al-awwal.}
\end{quote}

\begin{quote}
The bridge was built in the reign of the first king.
\end{quote}

Plural: \ar{عهود}, \emph{ʿuhūd}.

\ar{في} \ar{عهد}: in the days of, under. Also a promise: \ar{وفى} \ar{بالعهد}, he kept the pledge.

\begin{cousinweb}
\textbf{\ar{عهد}} (\emph{ʿahd}) is built on the native Arabic root \emph{ʿ-h-d} (Lane), \textquotedblleft{}to enjoin, to bind by a compact\textquotedblright{}. A ruler's pledge marks his time, so \emph{ʿahd} is also a reign.
\end{cousinweb}

\begin{grammarlens}[title={What you've built}]
One more word for telling what happened.
\end{grammarlens}

\subsection*{Your turn}
\begin{itemize}
\raggedright
  \item \emph{Say it:} \emph{ʿahd}
  \item \emph{Say it:} \emph{ʿahd}, once more
  \item \emph{Say it:} \emph{buniya al-jisr fī ʿahd al-malik al-awwal.} — The bridge was built in the reign of the first king.
  \item \emph{Recall:} say the Arabic for an era, then the Arabic for at that time, then say \emph{ʿahd} again
\end{itemize}

\begin{tcolorbox}[breakable,colback=teal!4,colframe=teal!35!black,title={Before you move on}]
What does \ar{عهد} mean? (\textquotedblleft{}A reign, an era; a pledge\textquotedblright{}.) Say it once more.
\end{tcolorbox}
